\documentclass[12pt]{article}

\usepackage[margin=1.12in]{geometry}
\usepackage{amsmath,amssymb,amsthm,mathtools}
\usepackage{booktabs}
\usepackage{longtable}
\usepackage{hyperref}

\newtheorem{theorem}{Theorem}
\newtheorem{corollary}{Corollary}
\theoremstyle{definition}
\newtheorem{definition}{Definition}
\theoremstyle{remark}
\newtheorem{remark}{Remark}

\newcommand{\C}{\mathbb C}
\newcommand{\D}{\mathcal D}
\newcommand{\F}{\mathcal F}
\newcommand{\R}{\mathbb R}
\newcommand{\Z}{\mathcal Z}

\setlength{\emergencystretch}{3em}
\hypersetup{
  pdftitle={Finite-Center Accountability and the Positive Route to Binary Rank},
  pdfsubject={Operator algebras, finite centers, sufficiency, and AQFT reconstruction},
  pdfkeywords={von Neumann algebras, finite centers, Blackwell sufficiency, Petz sufficiency, AQFT, binary rank}
}

\title{Finite-Center Accountability and the Positive Route to Binary Rank}
\author{}
\date{Live review note of June 20, 2026}

\begin{document}
\maketitle

\begin{abstract}
Two finite-center obstruction notes show that quotient-visible operational,
sector, modular, recovery, continuum, or entropy data do not by themselves
reconstruct hidden finite central rank or the predicate \(q=2\).  This short
note states the matching positive criterion.  For a supplied finite center
\(\Z_F\simeq \ell^\infty(F)\), let \(V_F=\Z_F/\C 1\) be its contrast space.
A declared family of probes, effects, outcome statistics, sector labels,
recovery tasks, modular records, finite-alphabet laws, or direct measurements
generates a readout map \(R:V_F\to W\).  The supplied data determine exactly
the quotient \(V_F/\ker R\).  Hence full finite central cardinality is
recoverable only when \(R\) is faithful on the physical finite-center contrast
space and an admissibility or maximality premise says this represented center
is the whole physical finite center.  Under those premises \(q=\operatorname
{rank}R+1\).  A binary conclusion \(q=2\) additionally needs one accounted
contrast plus a theorem or measurement identifying that contrast as the full
finite center rather than a quotient of a larger hidden extension.

This is deliberately demotion-ready.  In finite statistical language it is a
Blackwell/sufficiency statement.  In operator-algebraic language it is a
Petz-style relative-to-the-supplied-channel boundary.  The possible value is
only the translation into black-hole/AQFT reconstruction: recovery, entropy,
sector, modular, or locally covariant records force \(q\) or \(q=2\) only
after finite-center accountability has been supplied.
\end{abstract}

\paragraph{Status.}
This is live non-frozen review metadata.  It does not modify either frozen Team
1 paper, Lean archive, public mirror manifest, or SHA-256 package.  Its purpose
is external classification: known theorem, false statement, missing hypothesis,
too broad physically, or useful expository corollary.

\section{Question}

Positive reconstruction theorems reconstruct the object supplied to them: a
represented von Neumann algebra, a locally covariant functor, a sector
category, a correctable code algebra, a channel tested by specified
preparations and effects, or a Type II algebra with trace and normalization
data.  The finite-center obstruction is the converse question.  If the
displayed data have already passed through a quotient or coarse readout, what
extra premise is needed before hidden finite central rank, and in particular
the binary predicate \(q=2\), becomes reconstructible?

\section{Finite Accountability}

\begin{definition}[finite-center contrast and readout]
Let \(F\) be a finite set with \(q=|F|\), and let
\[
  \Z_F=\ell^\infty(F),\qquad V_F=\Z_F/\C 1 .
\]
Thus \(\dim V_F=q-1\).  A finite-center readout is a linear map
\[
  R:V_F\longrightarrow W
\]
generated by the declared observation family: normal probes, effects, outcome
statistics, sector records, recovery tasks, modular records, a finite-alphabet
law, or a direct measurement channel.  Its null directions are
\[
  N_R=\ker R .
\]
\end{definition}

\begin{theorem}[finite accountability quotient]
\label{thm:accountability}
Let \(D_R\) be any data record whose dependence on finite-center contrasts
factors through the readout map \(R:V_F\to W\).  Then:
\begin{enumerate}
  \item \(D_R\) factors through \(V_F/N_R\).
  \item No invariant of the finite center that varies inside a coset of
    \(N_R\) is reconstructible from \(D_R\).
  \item The full finite central cardinality \(q\) is reconstructible from
    \(D_R\) only after adding both \(N_R=0\) and an admissibility or maximality
    premise saying that the represented finite center is the whole physical
    finite center, not a quotient, null sector, gauge reduction, excluded
    sector, or hidden extension.
  \item Under those added premises,
    \[
      q=\dim(V_F/N_R)+1=\operatorname{rank}R+1 .
    \]
\end{enumerate}
\end{theorem}

\begin{proof}
If two contrasts \(v,v'\in V_F\) have the same image under \(R\), then
\(v-v'\in N_R\).  Every record built from \(R\) assigns the same value to
\(v\) and \(v'\), so it factors through \(V_F/N_R\).  Consequently any target
that changes within a coset of \(N_R\) is not a function of the supplied data.

The finite commutative algebra \(\ell^\infty(F)\) has one atomic coordinate
for each element of \(F\).  After constants are quotiented out, its contrast
space has dimension \(q-1\).  If \(N_R=0\), the readout is faithful on that
contrast space, so it determines \(\dim V_F\).  If the admissibility premise
also says the represented finite center is physically exhaustive, this
dimension is the physical finite-center contrast dimension, and
\(q=\dim V_F+1\).  If \(N_R\neq0\), or if the represented center is only a
quotient or admissible subcenter of a larger hidden center, refinements inside
the null or hidden directions change the unaccounted finite cardinality while
preserving all \(R\)-visible data.
\end{proof}

\begin{corollary}[binary rank requires two inputs]
\label{cor:binary}
A conclusion \(q=2\) from finite-center readout data requires both:
\[
  \dim(V_F/N_R)=1
\]
and a maximality, admissibility, direct-measurement, or finite-alphabet theorem
identifying that accounted one-dimensional contrast as the full physical
finite center.  One accounted contrast alone proves only that the declared
experiment sees a binary quotient.
\end{corollary}

\section{Relation to Standard Sufficiency}

The theorem should be demoted if a standard source already states this exact
translation.  Blackwell comparison of experiments makes sufficiency relative
to a declared experiment or decision problem \cite{Blackwell1951,
Blackwell1953}.  Petz sufficiency makes recovery and equality of relative
entropy relative to a declared subalgebra or channel \cite{Petz1986,Petz1988}.
Von Neumann algebra ideal theory says ultraweakly closed two-sided ideals are
central-projection summands \cite{LurieVN,HoudayerVN}.  None of these facts is
new.

The possible finite-center contribution is the reconstruction boundary:
black-hole/AQFT records determine a finite central cardinality only when their
hypotheses include faithful finite-center accountability.  Without that
premise, quotient-visible entropy, recovery, sector, modular, or locally
covariant records reconstruct the supplied quotient, not hidden cardinality.

\section{Claim Boundary for Reviewers}

\begin{longtable}{p{0.36\linewidth}p{0.52\linewidth}}
\toprule
Claim & Status \\
\midrule
\(R\) determines \(V_F/\ker R\). &
Elementary finite-dimensional linear algebra. \\
\addlinespace
\(R\) determines full \(q\) when \(\ker R=0\) and the represented center is
physically exhaustive. &
Positive accountability criterion. \\
\addlinespace
One accounted central contrast implies \(q=2\). &
False without maximality, admissibility, direct measurement, or a
finite-alphabet theorem. \\
\addlinespace
Blackwell/Petz sufficiency may already contain the criterion. &
Reviewer should identify the exact theorem if so; Team 1 should demote. \\
\addlinespace
AQFT/OAQEC hypotheses may already supply maximality or faithful support. &
Reviewer should name the hypothesis and decide whether it is present in the
Team 1 statement. \\
\bottomrule
\end{longtable}

\section{Reviewer Question}

Is Theorem~\ref{thm:accountability} merely finite Blackwell/Petz sufficiency
under a different name, or does the black-hole/AQFT translation isolate a useful
claim boundary?  If it is standard, the desired response is the citation and
the exact theorem name.  If it is false or too broad, the desired response is a
counterexample or missing physical hypothesis.

\begin{thebibliography}{99}

\bibitem{Blackwell1951}
D. Blackwell,
Comparison of experiments,
in \emph{Proceedings of the Second Berkeley Symposium on Mathematical
Statistics and Probability}, University of California Press, 1951, 93--102.

\bibitem{Blackwell1953}
D. Blackwell,
Equivalent comparisons of experiments,
\emph{Annals of Mathematical Statistics} \textbf{24} (1953), 265--272.

\bibitem{Petz1986}
D. Petz,
Sufficient subalgebras and the relative entropy of states of a von Neumann
algebra,
\emph{Communications in Mathematical Physics} \textbf{105} (1986), 123--131.

\bibitem{Petz1988}
D. Petz,
Sufficiency of channels over von Neumann algebras,
\emph{Quarterly Journal of Mathematics} \textbf{39} (1988), 97--108.

\bibitem{LurieVN}
J. Lurie,
\emph{Math 261y: von Neumann Algebras}, Lecture 9, Proposition 1.

\bibitem{HoudayerVN}
C. Houdayer,
\emph{An Invitation to von Neumann Algebras}, lecture notes, Exercise 3.7.

\bibitem{BFV}
R. Brunetti, K. Fredenhagen, and R. Verch,
The generally covariant locality principle: a new paradigm for local quantum
field theory,
\emph{Communications in Mathematical Physics} \textbf{237} (2003), 31--68.

\bibitem{FewsterVerch}
C. J. Fewster and R. Verch,
Dynamical locality and covariance: what makes a physical theory the same in
all spacetimes?,
\emph{Annales Henri Poincare} \textbf{13} (2012), 1613--1674.

\bibitem{JLMS}
D. L. Jafferis, A. Lewkowycz, J. Maldacena, and S. J. Suh,
Relative entropy equals bulk relative entropy,
\emph{Journal of High Energy Physics} \textbf{2016}, 004 (2016).

\end{thebibliography}

\end{document}
